\section{Assertion Contract}

A full RAIRE/AWAIRE audit record must make the implication from assertion set to
IRV winner public. RCOUNT currently records the first layer of that contract:
which ranked-choice method is being claimed, which candidates appear in the
reported elimination order, which assertions the package declares, and which
ranked choices support each sampled assertion value.

\begin{center}
\small
\begin{tabularx}{\textwidth}{lX}
\toprule
Field & Meaning \\
\midrule
\texttt{method\_id} & Versioned ranked-choice audit method. Current ids are \texttt{raire-irv-v1} and \texttt{awaire-irv-v1}. \\
\texttt{rcv\_elimination\_order} & Reported IRV elimination order, with the reported winner last in the synthetic boundary fixture. \\
\texttt{assertions} & Declared ranked-choice assertions. The current kind is \texttt{ranked-choice-assertion}. \\
\texttt{assorter\_id} & Assertion family or replay id, such as an NEB-style RAIRE id or an AWAIRE adaptive id. \\
\texttt{sample\_steps} & Sampled ranked ballots or CVRs tied to an assertion id and an assorter value. \\
\texttt{ranked\_choices} & Ordered candidate choices observed for the sampled unit. \\
\bottomrule
\end{tabularx}
\end{center}

The intended mathematical contract is stronger than the implemented one. A
RAIRE record should identify NEN and NEB assertions and show that their joint
truth implies the reported IRV winner. An AWAIRE record should additionally
make the adaptive weighting and assertion-risk path replayable. V.19 records
those claims as explicit boundary evidence, but it does not yet prove the
IRV-winner implication or recompute assertion statistics.
